\documentclass[10pt,a4paper]{amsart}
\usepackage[margin=19mm]{geometry}
\usepackage{amsmath,amssymb,mathtools}
\usepackage[scr=rsfs,cal=euler]{mathalfa}
\usepackage{xfrac}
\usepackage[unicode,hidelinks,bookmarks=false]{hyperref}
\newcommand{\ie}{i.e.\ }
\newcommand{\cf}{cf.\ }
\newcommand{\resp}{respectively\ }
\newcommand{\Subsection}{\subsection}
\providecommand{\nobreakdash}{\nobreak\hbox{-}}
\setlength{\emergencystretch}{2em}
\multlinegap0pt
\usepackage{amssymb}
\usepackage{dsfont}
\usepackage{tikz}
\newtheorem{lemma}{Lemma}
\title{A Correspondence between Billiards and Geodesics}
\author{Daniele Giannetto}
\date{Source: arXiv:2602.02938v3 (CC BY 4.0)}
\begin{document}
\maketitle

\noindent\emph{Source and licence.} A Correspondence between Billiards and Geodesics. Version arXiv:2602.02938v3 (CC BY 4.0), distributed by arXiv under the Creative Commons Attribution 4.0 International licence, http://creativecommons.org/licenses/by/4.0/, as recorded in the arXiv licence metadata for this exact version and shown on the abstract page https://arxiv.org/abs/2602.02938v3. Full licence notice: \texttt{licenses/arxiv-2602.02938-CC-BY-4.0.txt}.\par\smallskip

\noindent\textit{Editorial scope.} Counted unit: the article's lemma lem:coord\_connectors (source0.tex lines 903--908). The article's complete original proof of it is delivered with it (source0.tex lines 910--940, one \texttt{proof} environment of 31 lines). No mathematics is written, repaired or re-derived: the statement and the proof are the article's own lines, in the article's own notation, and no retained line is substituted or re-flowed.  No further source-local context is needed: every cross-reference of every retained line resolves inside the two delivered spans.  Nothing was omitted from any retained span, so the package carries no omission record.  No retained line contains a citation, so the wrapper carries no bibliography and no bibliography entry.  No retained line contains a figure environment, an image-inclusion command or drawing code, so no image is delivered and the package declares no asset dependency.  Editorial formatting only, in addition: an AMS \texttt{amsart} preamble that restates the article's own declarations and macros used by the retained lines, namely the environments \texttt{lemma} on separate counters, together with the standard packages those lines need.  The retained environments print in this document as Lemma 1 in order of appearance, whereas the article prints the same items with the numbering of its own full document.  The complete native source of the article as distributed in the arXiv e-print for this exact version is bundled unchanged as \texttt{sources/193873/source0.tex}.
\begin{lemma}\label{lem:coord_connectors}
In the context of the proof of the above proposition, for every $\eta>0$ there exists
$B_\eta<\infty$ depending only on $\eta$ and $C$ such that, for every $q_1,q_2\in C$, there is a Lipschitz curve $\alpha:[0,1]\to C$ joining $q_1$ to $q_2$ and satisfying
$$L_{D(K)}(\alpha)\le d_{D(K)}(q_1,q_2)+\eta,\qquad
\operatorname{Var}(\sigma\circ\alpha)\le B_\eta.$$
\end{lemma}

\begin{proof}
Fix $\eta>0$ and let $C_1\subset D(K)$ be compact and satisfy $C\subset C_1$, $C\cap \partial C_1 = \varnothing$. Since $C$ is compact, we may cover $C$ by finitely many coordinate charts $V_1,\ldots,V_m\Subset C_1$ with the following
property: if two points $q_1,q_2\in C$ are sufficiently close and lie in one of these $V_j$'s, then they can be connected inside that
neighborhood by a coordinate connector $\beta$ (i.e. a curve that, in coordinates, is a straight path) satisfying
$$L_{D(K)}(\beta)<\frac{\eta}{8},
\qquad
\operatorname{Var}(\sigma\circ\beta)\le A,$$
where $A<\infty$ is independent of $q_1,q_2$.
If $V_j$ is a coordinate chart away from $\partial K$, the property follows from the fact that $\sigma$ is smooth in $V_j$. Around some $q_0\in \partial K$, we use local topological coordinates $(\theta,\sigma)$ (where $\theta$ is a smooth coordinate for $\partial K$ around $q_0$ and $\sigma$ is a continuous but not differentiable coordinate) to construct local connectors of the form
$$(\theta_1+ t(\theta_2-\theta_1),\sigma_1+t(\sigma_2-\sigma_1)),$$
which clearly have finite $\sigma$-variation. We can also take the coordinate charts small enough so that the length of those coordinate connectors is less than $\eta/8$.

Now, we choose $r<\eta/8$ small enough so that every two points of $C$ with $d_{D(K)}$-distance less than $r$ can be joined by one of the above local connectors. Also, since $C$ is compact, we consider $\{a_1,\ldots,a_N\}\subset C$ such that for each $q\in C$ there is $j\in \{1,\ldots,N\}$ such that $d_{D(K)}(q,a_j)<r$. For each pair $a_i,a_j$, fix a Lipschitz curve $\gamma_{ij}:[0,1]\to C_1$ joining $a_i$ to $a_j$ such that
$$L_{D(K)}(\gamma_{ij})
\le d_{D(K)}(a_i,a_j)+\frac{\eta}{2},
\qquad
\operatorname{Var}(\sigma\circ\gamma_{ij})<\infty.$$
Since there are only finitely many pairs, we define $B_0:=\max\limits_{1\le i,j\le N} \operatorname{Var}(\sigma\circ\gamma_{ij})<\infty$.

Now let $q_1,q_2\in C$ and choose $a_i,a_j$ such that $d_{D(K)}(q_1,a_i)$, $d_{D(K)}(q_2,a_j)<r$. Let $\beta_1$ be a local connector from $q_1$ to $a_i$, and let $\beta_2$ be a local connector from $a_j$ to $q_2$. Define $\alpha$ as the concatenation of $\beta_1$, $\gamma_{ij}$ and $\beta_2$ in this order. We estimate
$$\operatorname{Var}(\sigma\circ\alpha) \le \operatorname{Var}(\sigma\circ\beta_1) + \operatorname{Var}(\sigma\circ\gamma_{ij}) + \operatorname{Var}(\sigma\circ\beta_2) \le 2A+B_0.$$
Thus we set $B_\eta:=2A+B_0$. 
To conclude that $\alpha$ is one of the curves we were looking for, we observe
\begin{equation*}
\begin{split}
L_{D(K)}(\alpha) &\le L_{D(K)}(\beta_1) + L_{D(K)}(\gamma_{ij}) + L_{D(K)}(\beta_2)\\
&\le \frac{\eta}{8} + d_{D(K)}(a_i,a_j) + \frac{\eta}{2} + \frac{\eta}{8}\\
&\le d_{D(K)}(q_1,q_2) +\frac{\eta}{4}+\frac{3\eta}{4} = d_{D(K)}(q_1,q_2) +\eta.
\end{split}
\end{equation*}
\end{proof}

\end{document}
